\section{Population-scoped discrepancy allowance}
\label{sec:b-theory}

The measurement precondition is an explicitly sourced allowance for
outcome-relevant reconstruction error, not sampling uncertainty.  Let $D$ be a
declared \emph{joint discrepancy event} covering the relevant decision
information and both regime outcomes, with $\Prb(D)\le p$.  Utility lies in
$[0,1]$, so every signed contrast $d(W)\,(U^v-U^0)$ lies in $[-1,1]$.

\begin{proposition}[population bound; application of the companion's Theorem]
\label{prop:n1-population}
Under the declared joint discrepancy event of probability at most $p$, bounded
utility, a \emph{fixed} policy $d$, and the companion's calibration
independence, a population lower confidence bound on the intended contrast loses
at most $2p$: if $L$ is a valid lower bound for the no-discrepancy population
gain $G(d)=\E[d(W)(U^v-U^0)]$, then $L-2p$ is a valid lower bound for the
population gain under the declared class.
\end{proposition}

\begin{proof}[Proof sketch]
The contrast is unchanged off $D$.  On $D$ the integrand lies in $[-1,1]$, so
moving any mass at most $p$ between the two regime outcomes changes the
population mean by at most $2p$.  This is the $2p$ perturbation bound of the
anonymous companion's bounded-outcome corollary \citep{companion2026}, applied
to a fixed policy at population scope; it is not a new safety theorem.
\end{proof}

\paragraph{Scope: no pointwise gate.}
A global $p$ supports a \emph{population} statement for a fixed policy only.  It
does \emph{not} justify a pointwise gate $\mathrm{LCB}(w)-2p$ applied per
checkpoint, nor a conditional clean-stratum guarantee without a separately
justified conditional allowance; selection can concentrate every defect.  The
binary $R=1$ sharper endpoint of the companion does not transfer to continuous
$U$ without a new derivation and is not used here.

\paragraph{Finite frozen families.}
Selecting among a finite frozen policy family uses \emph{simultaneous} population
lower bounds (a stated Bonferroni or equivalent allocation), and the procedure
abstains if no member qualifies.  Selection never proceeds by the pointwise rule
above.

\paragraph{Path allowance and allocation interference.}
For the two-checkpoint comparison the allowance is $p_{\text{path}}$, covering
the decision histories and the compared regime outcomes.  When only
checkpoint-wise bounds $p_1,p_2$ are available, they combine by the justified
union bound $\min(1,p_1+p_2)$, not by independence; a one-step allowance does not
automatically qualify a two-step trajectory.  Because all controllers share one
batch allocator, a wrong score at one root can change another root's allocation.
The population tax is therefore $2\,m\,p_{\text{path}}$, where $m$ is the maximum
number of roots whose realized regime can change when one root's recorded
information changes.  For the two-phase hard-budget allocator we derive $m\le 3$
(the root itself, at most one root displaced in phase 1, and at most one in phase
2 through the capacity shift); a single-root perturbation check attains 3.
\textbf{N1 qualifies conclusions; it does not produce a measurement-aware
controller or a deployment safety guarantee.}  In the sensitivity analysis each
grid value $p$ is a per-checkpoint allowance ($p_1=p_2=p$), so
$p_{\text{path}}=\min(1,2p)$; at the working value $p=0.02$ the population tax is
$2\cdot 3\cdot 0.04=0.24$ in $U$.  The same mapping is used for the grid, the
tax and the reported bound.  The p-grid report states that it covers allocation
interference under the derived $m\le 3$; the bound was reproduced by an
independent AI-assisted review, which is not a human proof review.

\paragraph{Provenance and sensitivity.}
The allowance is a \emph{sensitivity assumption} on the registered grid
$\{0,\,0.005,\,0.01,\,0.02,\,0.05,\,0.10\}$ with working value $0.02$; it is
never called measured.  The evaluation reports the largest $p$ at which each
frozen conclusion survives, adjusted if that largest $p$ is itself searched.
The contract records
\texttt{replay\_allowance.provenance\_kind}$\in\{\text{sensitivity\_assumption},\allowbreak
\text{validated\_evidence}\}$ and keeps \texttt{measured\_p} null until evidence
exists.  Task-condition prevalence is renamed $\omega$ to avoid collision with
the replay allowance $p$.

\paragraph{Equivalence and displaced theory.}
The ridge advantage/action-value identity is an implementation equivalence check
(\cref{sec:b-equivalence}), and the nonlinear direct-contrast versus two-head
comparison is at most a matched secondary ablation.  The inherited ordinal,
abstraction and price-theoretic material that is not needed for these claims is
preserved in \texttt{papers/paper\_b/archive/B1\_20260927/} and is not part of
the main narrative or the claim ledger.
